\\documentclass[11pt]{article}
\\usepackage[T1]{fontenc}
\\usepackage[utf8]{inputenc}
\\usepackage{geometry}
\\usepackage{xcolor}
\\usepackage{fvextra}
\\usepackage{hyperref}
\\geometry{margin=1in}
\\DefineVerbatimEnvironment{CodeBlock}{Verbatim}{breaklines,breakanywhere,fontsize=\\small}

\begin{document}
\section*{core/phase\\_space\\_integrator.py}
\textbf{Source Path}: \texttt{core/phase\\_space\\_integrator.py}\par
\bigskip
\begin{CodeBlock}
"""
phase_space_integrator.py — The Final Exit Gate (Post-Track Assembly)

EXECUTIVE MANDATE: THE PHASE-SPACE INTEGRATOR (V1.2)

This module performs final assembly of all tracks into a single IntegratedParticle
that lives on a Unit Hypersphere. Distance on this sphere is the Absolute Metric
of Framing Difference.

THE SEVEN FINGERS (ASTER v3.2):
===============================
Track 1 (Logits):       The "Surface Claim" — explicit NLI verdict
Track 1.5 (Spectral):   The "Internal Stress" — spectral polarity gradient (∇Φ)
Track 2 (Hologram):     The "Base Terrain" — RKS-projected embeddings (δ_μν)
Track 3 (Blinker):      The "Conformal Factor" — Dirichlet variance (1/ρ)
Track 4 (Walker):       The "Kinetic Probe" — MCMC path conductivity
Track 5 (Hadamard):     The "Metric Assembly" — K_final = K_T2 ∘ K_T1.5 / sqrt(ρ⊗ρ)
Track 6 (HoTT):         The "Formal Proofs" — type-theoretic verdicts (✓/✗/?)

EXIT GATE (THIS MODULE): Final normalization to unit hypersphere after Track 5

THE PHYSICS:
============
- High Blinker + Low Walker Resistance = NOISE (poorly written, incoherent)
- High Blinker + High Walker Resistance = STRUCTURAL SINGULARITY (model fighting
  to maintain specific framing despite ambiguity — the "Hidden Crack")
- Low Blinker + High Walker Resistance = IDEOLOGICAL BARRIER (clear but hard to cross)
- Low Blinker + Low Walker Resistance = CONSENSUS (clear and uncontested)

THE METRIC TENSOR (ASTER v3.2):
===============================
When use_hadamard_fusion=True, Track 5 (hadamard_fusion.py) assembles:
    g_μν = (1/ρ)·δ_μν + ∇_μΦ·∇_νΦ

This Exit Gate (Track 6) then:
1. Receives fused hologram from Track 5 (Hadamard kernel → spectral embedding)
2. Concatenates with Track 1 (Logits) and Track 4 (Walker)
3. Applies unit-sphere projection for absolute distance metric

Author: Belief Transformer Project
"""

from __future__ import annotations

import torch
import torch.nn.functional as F
import numpy as np
from dataclasses import dataclass, field
from typing import Dict, List, Optional, Tuple, Any
from pathlib import Path

from .pca_removal import fit_whitening_matrix, apply_whitening_fixed
from .thermo_config import ThermodynamicConfig


# =============================================================================
# INTEGRATED PARTICLE: The Final Output
# =============================================================================

@dataclass
class IntegratedParticle:
    """
    The unified representation of an article in Phase Space.

    Every article becomes a point on the Unit Hypersphere where distance
    is the Absolute Metric of Framing Difference.
    """
    # Core vector (L2-normalized, lives on unit hypersphere)
    vector: torch.Tensor  # [D_total] where D_total = sum of all track dims

    # Per-track contributions (before concatenation, after scaling)
    track_contributions: Dict[str, torch.Tensor] = field(default_factory=dict)

    # Liar Score: disagreement between Surface (T1) and Stress (T1.5)
    # High liar score = article's explicit claim contradicts its internal tension
    liar_score: float = 0.0

    # Physics interpretation
    blinker_variance: float = 0.0   # Track 3: static crack magnitude
    walker_resistance: float = 0.0  # Track 4: kinetic crack magnitude

    # Derived classification
    singularity_type: str = "unknown"  # noise | structural | barrier | consensus

    # Terrain State (ASTER v3.2 Metric Geometry)
    # Based on g_μν = δ_μν + ∇Φ ⊗ ∇Φ where:
    #   - Density (δ): Article coverage / local neighborhood
    #   - Stress (∇Φ): Observer gradient / conflict magnitude
    # BRIDGE: High density + Low stress → smooth conductive path
    # SWAMP: High density + High stress → viscous turbulent region
    # TIGHTROPE: Low density + Low stress → narrow consensus corridor
    # VOID: Low density + High stress → metric singularity (rupture zone)
    terrain_state: str = "unknown"  # BRIDGE | SWAMP | TIGHTROPE | VOID

    # Phantom Path Differential (ASTER v3.2)
    # Verdicts: TAUTOLOGY | HONEST | PHANTOM | RUPTURE
    phantom_verdict: str = "unknown"
    phantom_delta: float = 0.0        # Divergence ratio δ = W/d
    d_spectral: float = 0.0           # Weighted spectral distance (expected cost)
    w_actual: float = 0.0             # Actual manifold work (Walker integral)
    phantom_confidence: float = 0.0   # Confidence in the verdict [0, 1]
    walker_state: str = "unknown"     # Walker terminal state

    # Provenance
    article_id: Optional[str] = None

    # Hadamard Fusion diagnostics (ASTER v3.2)
    hadamard_diagnostics: Optional[Dict[str, Any]] = None

    def to_dict(self) -> Dict[str, Any]:
        """Serialize for JSON/storage."""
        return {
            'vector': self.vector.cpu().numpy().tolist(),
            'track_contributions': {
                k: v.cpu().numpy().tolist()
                for k, v in self.track_contributions.items()
            },
            'liar_score': self.liar_score,
            'blinker_variance': self.blinker_variance,
            'walker_resistance': self.walker_resistance,
            'singularity_type': self.singularity_type,
            'terrain_state': self.terrain_state,
            'article_id': self.article_id,
            # Phantom Path Differential (ASTER v3.2)
            'phantom_verdict': self.phantom_verdict,
            'phantom_delta': self.phantom_delta,
            'd_spectral': self.d_spectral,
            'w_actual': self.w_actual,
            'phantom_confidence': self.phantom_confidence,
            'walker_state': self.walker_state,
            'hadamard_diagnostics': self.hadamard_diagnostics,
            'u_axis': self.track_contributions.get('antagonism', torch.zeros(1)).cpu().numpy().tolist(),
        }


@dataclass
class IntegratorConfig:
    """Configuration for PhaseSpaceIntegrator."""

    # Robust scaling sample size
    scaling_sample_size: int = 5000

    # ZCA regularization
    zca_regularization: float = 1e-5

    # Track dimension expectations (for validation)
    expected_dims: Dict[str, int] = field(default_factory=lambda: {
        'logits': 24,           # [8 bots × 3 classes]
        'antagonism': 768,      # Gradient dimension (model hidden size)
        'hologram': 2048,       # RKS dim (configurable)
        'blinker': 2048,        # Dirichlet fused_std dim
        'walker': 2048,         # MCMC projected dim
    })

    # Physics thresholds for singularity classification
    blinker_high_threshold: float = 0.7   # High variance = blurry
    walker_high_threshold: float = 0.7    # High resistance = hard to traverse
    liar_threshold: float = 0.5           # Surface vs stress disagreement

    # Whether to apply ZCA to concatenated stack (ablation only — disabled by default)
    apply_global_zca: bool = False

    # Whether to project to unit sphere
    project_to_sphere: bool = True


# =============================================================================
# PHASE SPACE INTEGRATOR: The Fusion Engine
# =============================================================================

class PhaseSpaceIntegrator:
    """
    The Final Exit Gate: Fuses all six tracks into a single unit-sphere particle.

    Architecture:
    1. Receive all track tensors
    2. Apply robust scaling (fit on 5k sample, apply to all)
    3. Concatenate into mega-vector
    4. Apply global ZCA whitening (decorrelate tracks)
    5. L2 normalize to unit hypersphere
    6. Compute physics diagnostics (liar score, singularity type)
    """

    def __init__(self, config: Optional[IntegratorConfig] = None):
        self.config = config or IntegratorConfig()

        # Scaling parameters (fit once, apply to all)
        self.scaling_params: Dict[str, Dict[str, torch.Tensor]] = {}
        self.is_fit = False

        # Global ZCA parameters
        self.global_zca_mu: Optional[torch.Tensor] = None
        self.global_zca_W: Optional[torch.Tensor] = None

    def fit(
        self,
        track_samples: Dict[str, torch.Tensor],
    ) -> None:
        """
        Fit the integrator on a sample corpus (typically 5000 articles).

        This computes:
        1. Per-track robust scaling parameters (median, IQR)
        2. Global ZCA whitening matrix on concatenated tracks

        Args:
            track_samples: Dict mapping track names to [N, D] tensors
                Expected keys: 'logits', 'antagonism', 'hologram', 'blinker', 'walker'
        """
        device = next(iter(track_samples.values())).device

        # 1. Fit robust scaling for each track (with dimension validation)
        print("[PhaseSpaceIntegrator] Fitting robust scaling...")
        N_fit = None  # Will be set from first valid 2D tensor
        for track_name, samples in track_samples.items():
            if samples is None or samples.numel() == 0:
                continue

            # Ensure tensor is 2D [N, D]
            if samples.dim() == 1:
                # 1D tensor - skip for fitting (need batch dim)
                print(f"  Track '{track_name}': skipping 1D tensor {samples.shape} for fitting")
                continue
            elif samples.dim() == 3:
                # 3D tensor [N, B, D] - flatten to [N, B*D]
                samples = samples.reshape(samples.shape[0], -1)
                track_samples[track_name] = samples
                print(f"  Track '{track_name}': flattened 3D to {samples.shape}")
            elif samples.dim() != 2:
                print(f"  Track '{track_name}': skipping unexpected dim {samples.dim()}")
                continue

            if N_fit is None:
                N_fit = samples.shape[0]

            self.scaling_params[track_name] = self._fit_robust_scaler(samples)
            print(f"  Track '{track_name}': shape {samples.shape}, "
                  f"median_norm={self.scaling_params[track_name]['median_norm']:.4f}")

        # 2. Scale all tracks and concatenate (only 2D tensors with scaling params)
        scaled_tracks = []
        for track_name in ['logits', 'antagonism', 'hologram', 'blinker', 'walker']:
            if track_name in track_samples and track_samples[track_name] is not None:
                samples = track_samples[track_name]
                # Only include tracks that have scaling params (i.e., were valid 2D during step 1)
                if track_name not in self.scaling_params:
                    continue
                # Verify still 2D
                if samples.dim() != 2:
                    print(f"  [WARN] Track '{track_name}' not 2D after preprocessing, skipping")
                    continue
                scaled = self._apply_robust_scaling(samples, self.scaling_params[track_name])
                scaled_tracks.append(scaled)

        if not scaled_tracks:
            raise ValueError("No valid track samples provided for fitting")

        concatenated = torch.cat(scaled_tracks, dim=-1)  # [N, D_total]
        print(f"[PhaseSpaceIntegrator] Concatenated shape: {concatenated.shape}")

        # 3. Fit global ZCA on concatenated stack
        if self.config.apply_global_zca:
            print("[PhaseSpaceIntegrator] Fitting global ZCA whitening...")
            # Subsample if too large
            if concatenated.shape[0] > self.config.scaling_sample_size:
                idx = torch.randperm(concatenated.shape[0])[:self.config.scaling_sample_size]
                zca_sample = concatenated[idx]
            else:
                zca_sample = concatenated

            zca_fit = fit_whitening_matrix(zca_sample, self.config.zca_regularization)
            self.global_zca_mu = zca_fit['mu'].to(device)
            self.global_zca_W = zca_fit['whitening_matrix'].to(device)
            print(f"  ZCA fitted on {zca_sample.shape[0]} samples, "
                  f"condition={zca_fit['singular_values'][0] / zca_fit['singular_values'][-1]:.2f}")

        self.is_fit = True
        print("[PhaseSpaceIntegrator] Fitting complete.")

    def _fit_robust_scaler(self, samples: torch.Tensor) -> Dict[str, torch.Tensor]:
        """
        Compute robust scaling parameters (median and IQR-based scale).

        More robust to outliers than mean/std normalization.
        """
        # Compute per-sample L2 norms
        norms = samples.norm(dim=-1)  # [N]

        # Robust statistics
        median_norm = norms.median()
        q75 = torch.quantile(norms, 0.75)
        q25 = torch.quantile(norms, 0.25)
        iqr = q75 - q25

        # Scale factor: normalize so median norm becomes 1.0
        # Use IQR to handle spread
        scale = 1.0 / (median_norm + 1e-9)

        return {
            'median_norm': median_norm,
            'iqr': iqr,
            'scale': scale,
            'q25': q25,
            'q75': q75,
        }

    def _apply_robust_scaling(
        self,
        features: torch.Tensor,
        params: Dict[str, torch.Tensor],
    ) -> torch.Tensor:
        """Apply pre-computed robust scaling."""
        return features * params['scale'].to(features.device)

    def _apply_force_gain(
        self,
        position_tracks: Dict[str, torch.Tensor],
        force_tracks: Dict[str, torch.Tensor],
        min_energy_fraction: float = 0.2
    ) -> Dict[str, torch.Tensor]:
        """
        Ensure force tracks (antagonism, walker) have minimum energy representation.

        Without this, the high-dimensional hologram/blinker tracks dominate,
        making the physics signals invisible.
        """
        # Compute position track energy
        E_pos = sum(
            t.pow(2).sum() for t in position_tracks.values()
            if t is not None and t.numel() > 0
        )

        # Compute force track energy
        E_force = sum(
            t.pow(2).sum() for t in force_tracks.values()
            if t is not None and t.numel() > 0
        )

        if E_force < 1e-9 or E_pos < 1e-9:
            return force_tracks

        # Target: force should be at least min_energy_fraction of total
        # E_force / (E_pos + E_force) >= min_energy_fraction
        # E_force >= min_energy_fraction * E_pos / (1 - min_energy_fraction)
        target_E_force = min_energy_fraction * E_pos / (1 - min_energy_fraction)

        if E_force < target_E_force:
            gain = torch.sqrt(target_E_force / E_force)
            print(f"[PhaseSpace] Applying gain {gain:.2f} to force tracks (energy: {E_force:.4f} -> {target_E_force:.4f})")
            return {
                name: (t * gain if t is not None else None)
                for name, t in force_tracks.items()
            }

        return force_tracks

    def integrate(
        self,
        logits: Optional[torch.Tensor] = None,          # [N, 24] or [N, 8, 3]
        antagonism: Optional[torch.Tensor] = None,      # [N, 768]
        hologram: Optional[torch.Tensor] = None,        # [N, D_rks]
        blinker: Optional[torch.Tensor] = None,         # [N, D_rks] (fused_std)
        walker: Optional[torch.Tensor] = None,          # [N, D_rks] (MCMC projected)
        article_ids: Optional[List[str]] = None,
        use_hadamard_fusion: bool = False,              # ASTER v3.2: Hadamard kernel product
    ) -> List[IntegratedParticle]:
        """
        Integrate all tracks into IntegratedParticles.

        Args:
            logits: Track 1 - NLI verdict logits
            antagonism: Track 1.5 - Structural antagonism gradients
            hologram: Track 2 - ZCA-whitened Matern-projected embeddings
            blinker: Track 3 - Dirichlet variance (fused_std)
            walker: Track 4 - MCMC path conductivity projection
            article_ids: Optional article identifiers

        Returns:
            List of IntegratedParticle objects, one per article
        """
        if not self.is_fit:
            raise RuntimeError("PhaseSpaceIntegrator must be fit() before integrate()")

        # Determine batch size from first non-None tensor
        N = None
        for t in [logits, antagonism, hologram, blinker, walker]:
            if t is not None:
                N = t.shape[0]
                device = t.device
                break

        if N is None:
            return []

        # Flatten logits if needed [N, 8, 3] → [N, 24]
        if logits is not None and logits.dim() == 3:
            logits = logits.reshape(N, -1)

        # Collect and scale tracks
        tracks = {
            'logits': logits,
            'antagonism': antagonism,
            'hologram': hologram,
            'blinker': blinker,
            'walker': walker,
        }

        scaled_tracks = {}
        for name, tensor in tracks.items():
            if tensor is not None and name in self.scaling_params:
                scaled_tracks[name] = self._apply_robust_scaling(
                    tensor, self.scaling_params[name]
                )
            elif tensor is not None:
                # No scaling params — use as-is with warning
                print(f"[WARN] Track '{name}' has no scaling params, using raw values")
                scaled_tracks[name] = tensor

        # Apply gain to ensure force tracks are visible
        position_tracks = {k: scaled_tracks.get(k) for k in ['logits', 'hologram', 'blinker']}
        force_tracks = {k: scaled_tracks.get(k) for k in ['antagonism', 'walker']}
        boosted_force = self._apply_force_gain(position_tracks, force_tracks)
        scaled_tracks.update(boosted_force)

        # ASTER v3.2: Hadamard Kernel Product Fusion
        # Instead of concatenating Track 2 + Track 1.5 + Track 3, we:
        #   1. Hadamard product: K_final = K_rks * K_spectral (T2 * T1.5)
        #   2. Conformal scaling: distances scaled by Track 3 variance
        #   3. Spectral embedding: recover vector space for Walker
        hadamard_diagnostics = None
        if use_hadamard_fusion:
            thermo_config = ThermodynamicConfig()
            t2 = scaled_tracks.get('hologram')
            t15 = scaled_tracks.get('antagonism')
            t3 = scaled_tracks.get('blinker')

            if t2 is not None and t15 is not None:
                from .hadamard_fusion import compute_kernel_with_conformal, compute_hadamard_fusion_simple, HadamardFusionConfig

                print("[PhaseSpaceIntegrator] Applying Hadamard fusion (T2 * T1.5)...")

                hf_config = HadamardFusionConfig(
                    output_dim=t2.shape[-1],  # Match hologram dimension
                    dark_manifold_rescue=False,
                    hadamard_softening=0.0,
                    kernel_floor=thermo_config.hadamard_floor,
                )

                if t3 is not None:
                    # Full pipeline with conformal scaling from Track 3
                    hf_result = compute_kernel_with_conformal(t2, t15, t3, config=hf_config)
                    fused_hologram = hf_result.embeddings
                    hadamard_diagnostics = hf_result.to_dict()
                    print(f"  Hadamard fusion complete: {hf_result.n_isolated_rescued} nodes rescued, "
                          f"sparsity={hadamard_diagnostics['K_hadamard_sparsity']:.3f}")
                else:
                    # Simple Hadamard without conformal (no Track 3)
                    fused_hologram, hadamard_diagnostics = compute_hadamard_fusion_simple(t2, t15, config=hf_config)
                    print(f"  Hadamard fusion (no conformal): {hadamard_diagnostics['n_isolated_rescued']} nodes rescued")

                # Replace hologram with fused representation
                scaled_tracks['hologram'] = fused_hologram
                # Mark antagonism and blinker as consumed (not to be concatenated)
                scaled_tracks['antagonism'] = None
                scaled_tracks['blinker'] = None
            else:
                print("[WARN] Hadamard fusion requested but Track 2 or Track 1.5 missing, falling back to concatenation")

        # Concatenate in canonical order with dimension validation
        # Note: When use_hadamard_fusion=True, antagonism and blinker are None (already fused into hologram)
        concat_list = []
        track_slices = {}  # Track where each track lives in the concatenated vector
        offset = 0
        for name in ['logits', 'antagonism', 'hologram', 'blinker', 'walker']:
            if name in scaled_tracks and scaled_tracks[name] is not None:
                t = scaled_tracks[name]

                # Ensure tensor is 2D [N, D]
                if t.dim() == 1:
                    # 1D tensor [D] - expand to [N, D] by repeating
                    t = t.unsqueeze(0).expand(N, -1)
                    print(f"[WARN] Track '{name}' was 1D, expanded to {t.shape}")
                elif t.dim() == 3:
                    # 3D tensor [N, B, D] - flatten to [N, B*D]
                    t = t.reshape(t.shape[0], -1)
                    print(f"[WARN] Track '{name}' was 3D, flattened to {t.shape}")
                elif t.dim() != 2:
                    print(f"[WARN] Track '{name}' has unexpected dim {t.dim()}, skipping")
                    continue

                # Validate batch dimension matches
                if t.shape[0] != N:
                    print(f"[WARN] Track '{name}' batch size {t.shape[0]} != expected {N}, skipping")
                    continue

                concat_list.append(t)
                track_slices[name] = (offset, offset + t.shape[-1])
                offset += t.shape[-1]

        if not concat_list:
            return []

        concatenated = torch.cat(concat_list, dim=-1)  # [N, D_total]

        # Apply global ZCA whitening (with dimension check)
        if self.config.apply_global_zca and self.global_zca_W is not None:
            zca_dim = self.global_zca_W.shape[0]
            concat_dim = concatenated.shape[-1]
            if concat_dim != zca_dim:
                print(f"[WARN] ZCA dimension mismatch: concat {concat_dim} vs ZCA {zca_dim}, skipping ZCA")
            else:
                concatenated = apply_whitening_fixed(
                    concatenated,
                    self.global_zca_mu,
                    self.global_zca_W,
                )

        # Project to unit hypersphere
        if self.config.project_to_sphere:
            concatenated = F.normalize(concatenated, p=2, dim=-1)

        # Compute physics diagnostics and build particles
        # First pass: compute raw blinker/walker values for normalization
        raw_blinker = []
        raw_walker = []
        for i in range(N):
            if blinker is not None:
                raw_blinker.append(float(blinker[i].norm().item()))
            else:
                raw_blinker.append(0.0)

            if walker is not None:
                raw_walker.append(float(walker[i].norm().item()))
            else:
                raw_walker.append(0.0)

        # Normalize to [0, 1] using percentile scaling (robust to outliers)
        def percentile_normalize(values):
            """Normalize values to [0, 1] using min-max with percentile clipping."""
            arr = np.array(values)
            if arr.max() == arr.min():
                return np.zeros_like(arr)  # All same value
            # Use 5th-95th percentile to be robust to outliers
            p5, p95 = np.percentile(arr, [5, 95])
            if p95 == p5:
                p5, p95 = arr.min(), arr.max()
            clipped = np.clip(arr, p5, p95)
            normalized = (clipped - p5) / (p95 - p5 + 1e-10)
            return normalized

        norm_blinker = percentile_normalize(raw_blinker)
        norm_walker = percentile_normalize(raw_walker)

        # Second pass: build particles with normalized values
        particles = []
        for i in range(N):
            # Extract per-track contributions from the integrated vector
            track_contribs = {}
            for name, (start, end) in track_slices.items():
                track_contribs[name] = concatenated[i, start:end]

            # Compute Liar Score: disagreement between T1 (logits) and T1.5 (antagonism)
            liar_score = self._compute_liar_score(
                logits[i] if logits is not None else None,
                antagonism[i] if antagonism is not None else None,
            )

            # Use normalized blinker/walker for classification
            blinker_var = float(norm_blinker[i])
            walker_res = float(norm_walker[i])

            # Classify singularity type based on physics
            singularity = self._classify_singularity(
                blinker_var, walker_res, liar_score
            )

            # Classify terrain state (ASTER v3.2 metric geometry)
            terrain = self._classify_terrain(blinker_var, walker_res)

            particle = IntegratedParticle(
                vector=concatenated[i],
                track_contributions=track_contribs,
                liar_score=liar_score,
                blinker_variance=blinker_var,
                walker_resistance=walker_res,
                singularity_type=singularity,
                terrain_state=terrain,
                article_id=article_ids[i] if article_ids else None,
                hadamard_diagnostics=hadamard_diagnostics,  # ASTER v3.2
            )
            particles.append(particle)

        return particles

    def compute_phantom_differential(
        self,
        d_spectral: torch.Tensor,  # [N] — weighted spectral distance (expected cost)
        w_actual: List[float],     # [N] — actual manifold work (Walker integral)
        walker_states: List[str],  # [N] — Walker terminal states
        blinker_variance: Optional[List[float]] = None,  # [N] — Track 3 variance (for terrain)
    ) -> List[dict]:
        """
        Compute the Panic Function (Divergence Ratio) for Lie Detection.

        Delta = W_actual / d_spectral = (Manifold Work) / (Spectral Projection)

        Now also computes terrain_state (BRIDGE/SWAMP/TIGHTROPE/VOID) per article
        based on the density × stress manifold from Track 2.

        VERDICT TAXONOMY:
        =================
        TAUTOLOGY (δ < 0.8):
            Article claims controversy but embedding distance is tiny.
            "Much ado about nothing" — semantic noise, not real disagreement.

        HONEST (0.8 ≤ δ < 1.5):
            Manifold work matches spectral projection. The article's framing
            corresponds to genuine semantic distance. No hidden agenda.

        PHANTOM (1.5 ≤ δ < 10.0):
            Manifold work exceeds spectral expectation. The Walker had to do
            extra work navigating around semantic obstacles not visible in
            the flat projection. Hidden structure detected.

        RUPTURE (δ ≥ 10.0 or Walker broken/fog):
            Extreme divergence. The manifold path is completely disconnected
            from the spectral projection. Indicates severe framing distortion
            or fundamental semantic discontinuity.

        Args:
            d_spectral: [N] weighted spectral distance from SpectralPolarity
            w_actual: [N] actual work integrals from Walker (Track 4)
            walker_states: [N] terminal states from Walker ("elastic", "trapped", "broken", "fog")
            blinker_variance: [N] Track 3 variance values (optional, for terrain classification)

        Returns:
            List of verdict dicts with keys: verdict, delta, d_spectral, w_actual, confidence,
            walker_state, terrain_state, blinker_variance, walker_resistance
        """
        # Threshold constants (ASTER v3.2 calibration)
        TAUTOLOGY_THRESHOLD = 0.8   # δ < 0.8 → TAUTOLOGY
        HONEST_THRESHOLD = 1.5      # 0.8 ≤ δ < 1.5 → HONEST
        PHANTOM_THRESHOLD = 10.0    # 1.5 ≤ δ < 10.0 → PHANTOM, else RUPTURE

        # Minimum spectral distance to avoid division instability
        MIN_D_SPECTRAL = 0.01

        verdicts = []
        for i in range(len(w_actual)):
            d = float(d_spectral[i]) if d_spectral is not None else MIN_D_SPECTRAL
            w = w_actual[i]
            state = walker_states[i] if i < len(walker_states) else "elastic"

            # Get blinker variance for terrain classification
            blinker_var = blinker_variance[i] if blinker_variance and i < len(blinker_variance) else 0.5

            # Normalize walker resistance for terrain (high work = high resistance = low density)
            # Clamp to prevent infinity issues
            w_clamped = min(float(w) if w == w and w != float('inf') else 10.0, 100.0)
            walker_res = w_clamped / 10.0  # Normalize to ~[0, 10]

            # Classify terrain state (BRIDGE/SWAMP/TIGHTROPE/VOID)
            terrain = self._classify_terrain(blinker_var, walker_res)

            # Handle catastrophic Walker states first
            if state in ("fog", "broken") or w == float('inf') or w != w:  # w != w catches NaN
                verdicts.append({
                    "verdict": "RUPTURE",
                    "delta": float('inf'),
                    "d_spectral": d,
                    "w_actual": w if w == w else float('inf'),  # Replace NaN with inf
                    "confidence": 1.0,
                    "walker_state": state,
                    "terrain_state": terrain,
                    "blinker_variance": blinker_var,
                    "walker_resistance": walker_res,
                })
                continue

            # Handle zero/tiny spectral distance (no polarization detected)
            if d < MIN_D_SPECTRAL:
                # If Walker also did minimal work → HONEST (trivial case)
                # If Walker did significant work but no spectral distance → PHANTOM
                if w < MIN_D_SPECTRAL:
                    verdict = "HONEST"
                    delta = 1.0
                    confidence = 0.5  # Low confidence due to trivial case
                else:
                    # Spectral says "no difference" but Walker disagrees
                    verdict = "PHANTOM"
                    delta = w / MIN_D_SPECTRAL  # Approximate
                    confidence = 0.7
                verdicts.append({
                    "verdict": verdict,
                    "delta": delta,
                    "d_spectral": d,
                    "w_actual": w,
                    "confidence": confidence,
                    "walker_state": state,
                    "terrain_state": terrain,
                    "blinker_variance": blinker_var,
                    "walker_resistance": walker_res,
                })
                continue

            # Standard case: compute divergence ratio
            delta = w / d

            # Classify based on thresholds
            if delta < TAUTOLOGY_THRESHOLD:
                verdict = "TAUTOLOGY"
                # Confidence increases as delta approaches 0 (more clearly tautological)
                confidence = 1.0 - (delta / TAUTOLOGY_THRESHOLD)
            elif delta < HONEST_THRESHOLD:
                verdict = "HONEST"
                # Confidence peaks at delta=1.0 (perfect match)
                confidence = 1.0 - abs(delta - 1.0) / 0.5
            elif delta < PHANTOM_THRESHOLD:
                verdict = "PHANTOM"
                # Confidence increases with delta (more clearly phantom)
                confidence = min(1.0, (delta - HONEST_THRESHOLD) / (PHANTOM_THRESHOLD - HONEST_THRESHOLD))
            else:
                verdict = "RUPTURE"
                confidence = min(1.0, delta / PHANTOM_THRESHOLD - 1.0)

            verdicts.append({
                "verdict": verdict,
                "delta": delta,
                "d_spectral": d,
                "w_actual": w,
                "confidence": max(0.0, min(1.0, confidence)),
                "walker_state": state,
                "terrain_state": terrain,
                "blinker_variance": blinker_var,
                "walker_resistance": walker_res,
            })

        return verdicts

    def compute_divergence_ratio(
        self,
        spectral_result: Any,  # SpectralPolarityResult
        walker_work_integrals: torch.Tensor,  # [N]
        walker_states: List[str],
        emb_start: torch.Tensor,  # [N, H] — article embeddings at "start" pole
        emb_end: torch.Tensor,    # [N, H] — article embeddings at "end" pole
    ) -> List[dict]:
        """
        Full Panic Function computation using weighted spectral distance.

        This is the preferred entry point that uses the new signal-weighted
        spectral distance rather than the legacy unweighted version.

        Args:
            spectral_result: SpectralPolarityResult from SpectralPolarity.compute_batch()
            walker_work_integrals: [N] tensor of manifold work values from Walker
            walker_states: [N] list of Walker terminal states
            emb_start: [N, H] embeddings at one semantic pole
            emb_end: [N, H] embeddings at the other semantic pole

        Returns:
            List of verdict dicts (same format as compute_phantom_differential)
        """
        from .spectral_polarity import SpectralPolarity

        # Create SpectralPolarity instance for weighted distance computation
        sp = SpectralPolarity()

        # Compute weighted spectral distance using the spectral basis
        d_spectral = sp.compute_weighted_spectral_distance(
            emb_a=emb_start,
            emb_b=emb_end,
            singular_values=spectral_result.singular_values,
            Vt=spectral_result.u_basis,
            dynamic_k=spectral_result.dynamic_k,
        )

        # Convert walker work to list
        w_actual = walker_work_integrals.tolist() if isinstance(walker_work_integrals, torch.Tensor) else list(walker_work_integrals)

        # Use the main verdict computation
        return self.compute_phantom_differential(d_spectral, w_actual, walker_states)

    def _compute_liar_score(
        self,
        logits: Optional[torch.Tensor],
        antagonism: Optional[torch.Tensor],
    ) -> float:
        """
        Compute disagreement between Surface Claim (T1) and Internal Stress (T1.5).

        Liar Score = 1 - cos(direction_from_logits, direction_from_antagonism)

        High score = article's explicit verdict contradicts its internal tension.
        """
        if logits is None or antagonism is None:
            return 0.0

        # Convert logits to direction vector (which class is favored)
        if logits.dim() == 1 and logits.shape[0] == 24:
            # [24] → [8, 3] → mean over bots → [3]
            logits_3d = logits.reshape(8, 3).mean(dim=0)
        elif logits.dim() == 1 and logits.shape[0] == 3:
            logits_3d = logits
        else:
            logits_3d = logits.reshape(-1, 3).mean(dim=0) if logits.numel() >= 3 else logits

        # Softmax to get direction
        logits_dir = F.softmax(logits_3d, dim=-1)

        # For antagonism, we interpret the gradient direction
        # Normalize to unit vector
        antag_dir = F.normalize(antagonism.unsqueeze(0), p=2, dim=-1).squeeze(0)

        # Since they're different spaces, we use magnitude-based disagreement
        # High antagonism magnitude + confident logits but different "feel" = liar
        logits_confidence = float((logits_dir.max() - logits_dir.min()).item())
        antag_magnitude = float(antagonism.norm().item())

        # Liar score: high confidence + high stress = potential liar
        # Normalized to [0, 1]
        raw_score = logits_confidence * antag_magnitude
        liar_score = float(torch.sigmoid(torch.tensor(raw_score - 0.5)).item())

        return liar_score

    def _classify_singularity(
        self,
        blinker_var: float,
        walker_res: float,
        liar_score: float,
    ) -> str:
        """
        Classify the article's singularity type based on physics.

        The 2×2 matrix of Blinker × Walker:
        - High Blinker + Low Walker = NOISE (blurry and easy to traverse)
        - High Blinker + High Walker = STRUCTURAL (fighting to maintain frame)
        - Low Blinker + High Walker = BARRIER (clear but ideologically isolated)
        - Low Blinker + Low Walker = CONSENSUS (clear and mainstream)
        """
        b_high = blinker_var > self.config.blinker_high_threshold
        w_high = walker_res > self.config.walker_high_threshold

        if b_high and w_high:
            return "structural_singularity"
        elif b_high and not w_high:
            return "noise"
        elif not b_high and w_high:
            return "ideological_barrier"
        else:
            return "consensus"

    def _classify_terrain(
        self,
        blinker_var: float,
        walker_res: float,
        local_density: Optional[float] = None,
    ) -> str:
        """
        Classify the article's terrain state based on metric geometry.

        The 2×2 matrix of Density (δ) × Stress (∇Φ):

        Metric tensor: g_μν = δ_μν + ∇Φ ⊗ ∇Φ

        Variable A (δ): Article Density (The Floor)
            - High density = low walker resistance (easy manifold traversal)
            - Low density = high walker resistance (sparse/isolated)

        Variable B (∇Φ): Observer Gradient / Stress (The Stress Field)
            - High stress = high blinker variance (observer disagreement)
            - Low stress = low blinker variance (observer consensus)

        TERRAIN STATES:
        ===============
        BRIDGE: High Density + Low Stress
            → Smooth, stable, paved connection. Logic flows easily.

        SWAMP: High Density + High Stress
            → Ground exists but is muddy/turbulent. High viscosity.
            → Bots screaming at each other over a pile of data.

        TIGHTROPE: Low Density + Low Stress
            → Agreed upon but narrow. One step off and you fall.
            → Sparse but coherent corridor.

        VOID: Low Density + High Stress
            → Total collapse. The Metric Singularity.
            → Acts as an insulator, not empty space.
        """
        # Density: inverse of walker resistance (high resistance = low density)
        # Use threshold at 0.5 (normalized values)
        density_high = walker_res < self.config.walker_high_threshold

        # Stress: direct from blinker variance
        stress_high = blinker_var > self.config.blinker_high_threshold

        if density_high and not stress_high:
            return "BRIDGE"      # Solid + Consensus = Smooth path
        elif density_high and stress_high:
            return "SWAMP"       # Solid + Conflict = Viscous turbulence
        elif not density_high and not stress_high:
            return "TIGHTROPE"   # Sparse + Consensus = Narrow corridor
        else:
            return "VOID"        # Sparse + Conflict = Metric singularity

    def compute_riemannian_metric(
        self,
        features: torch.Tensor,
        blinker_variance: torch.Tensor,
        n_neighbors: int = 10,
    ) -> torch.Tensor:
        """
        Compute the Riemannian metric tensor contribution for each article.

        g_μν = δ_μν + ∇Φ ⊗ ∇Φ

        The stress term ∇Φ ⊗ ∇Φ acts as a viscosity field that
        increases geodesic cost in high-conflict regions.

        Args:
            features: [N, D] article embeddings
            blinker_variance: [N] normalized variance (stress proxy)
            n_neighbors: k for local density estimation

        Returns:
            [N] local metric coefficient (1 + stress²)
        """
        # Stress contribution: ||∇Φ||² ~ blinker_variance²
        stress_sq = blinker_variance ** 2

        # Metric coefficient: 1 + stress² (identity + stress tensor)
        metric_coeff = 1.0 + stress_sq

        return metric_coeff

    def save(self, path: Path) -> None:
        """Save integrator state for reproducibility."""
        state = {
            'config': self.config.__dict__,
            'scaling_params': {
                k: {kk: vv.cpu() for kk, vv in v.items()}
                for k, v in self.scaling_params.items()
            },
            'global_zca_mu': self.global_zca_mu.cpu() if self.global_zca_mu is not None else None,
            'global_zca_W': self.global_zca_W.cpu() if self.global_zca_W is not None else None,
            'is_fit': self.is_fit,
        }
        torch.save(state, path)
        print(f"[PhaseSpaceIntegrator] Saved to {path}")

    @classmethod
    def load(cls, path: Path) -> 'PhaseSpaceIntegrator':
        """Load integrator state."""
        state = torch.load(path, weights_only=False)
        config = IntegratorConfig(**state['config'])
        integrator = cls(config)
        integrator.scaling_params = {
            k: {kk: vv for kk, vv in v.items()}
            for k, v in state['scaling_params'].items()
        }
        integrator.global_zca_mu = state['global_zca_mu']
        integrator.global_zca_W = state['global_zca_W']
        integrator.is_fit = state['is_fit']
        print(f"[PhaseSpaceIntegrator] Loaded from {path}")
        return integrator


# =============================================================================
# CONVENIENCE FUNCTIONS
# =============================================================================

def create_integrator_from_pipeline_output(
    pipeline_outputs: List[Dict[str, Any]],
    config: Optional[IntegratorConfig] = None,
) -> PhaseSpaceIntegrator:
    """
    Create and fit a PhaseSpaceIntegrator from pipeline outputs.

    Args:
        pipeline_outputs: List of dicts from complete_pipeline.py runs
        config: Optional configuration

    Returns:
        Fitted PhaseSpaceIntegrator ready for integration
    """
    integrator = PhaseSpaceIntegrator(config)

    # Collect samples from all outputs
    track_samples = {
        'logits': [],
        'antagonism': [],
        'hologram': [],
        'blinker': [],
        'walker': [],
    }

    for output in pipeline_outputs:
        # Track 1: Logits
        if 'logits_raw' in output:
            track_samples['logits'].append(torch.tensor(output['logits_raw']))

        # Track 2: Hologram (fused embeddings)
        if 'dirichlet_fused' in output:
            track_samples['hologram'].append(torch.tensor(output['dirichlet_fused']))

        # Track 3: Blinker (variance)
        if 'dirichlet_fused_std' in output:
            track_samples['blinker'].append(torch.tensor(output['dirichlet_fused_std']))

    # Stack samples
    for name in track_samples:
        if track_samples[name]:
            track_samples[name] = torch.cat(track_samples[name], dim=0)
        else:
            track_samples[name] = None

    # Filter out None values
    track_samples = {k: v for k, v in track_samples.items() if v is not None}

    if track_samples:
        integrator.fit(track_samples)

    return integrator

\end{CodeBlock}
\end{document}
